\exercisesfor{7}

\begin{enumerate}
\def\labelenumi{\arabic{enumi}.}
\item
  设 \(\mathrm{G}\) 是李群, \(\phi: \mathrm{U} \to \mathrm{G}\) 是指数映射, 满足 \(\mathbf{Z}\mathbf{U} \subset \mathbf{U}\), 且对所有 \(\phi(rx) = \phi(x)^r\)\(x \in \mathbf{U}\) 和所有 \(r \in \mathbf{Z}\) 有 . 若 \(p > 0\), 则 \(\phi\) 是从 \(\mathrm{U}\) 到 \(\phi(\mathrm{U})\) 的解析同构. (若 \(\phi(x) = \phi(y)\), 则对所有  有 \(\phi(p^n x) = \phi(p^n y)\), 从而 \(n \in \mathbf{N}\)\(x = y\). 令 \(\mathrm{W}\) 为 \(0\)\(\mathrm{L}(\mathrm{G})\) 中 0 的一个开邻域, 使 \(\phi^{-1}\) 在 \(\phi(\mathrm{W})\) 上解析. 对所有 \(s \in \phi(\mathrm{U})\), 存在 \(n \in \mathbf{N}\) 及  中 s 的一个邻域 \(\mathrm{V}\), 使 \(s\)\(\phi(\mathrm{U})\)\(t \in \mathrm{V} \Rightarrow t^{p^n} \in \phi(\mathrm{W})\).)
\item
  设 G 是李群, \(\phi: \mathrm{U} \to \mathrm{G}\) 是具有命题 3 所述性质的指数映射, V 是 U 中 0 的一个邻域, 满足 \(\mathbf{ZV} \subset \mathrm{V}\), \(\psi: \mathrm{V} \to \mathrm{G}\) 是 \(\phi\) 在 0 处的切映射, 且对所有 \(\psi(nx) = \psi(x)^n\)\(n \in \mathrm{V}\) 和所有 \(n \in \mathbf{Z}\) 有 . 若 \(p > 0\), 则 \(\psi = \phi \mid \mathrm{V}\). (给 L(G) 赋予一个范数. 令 \(x \in \mathrm{V}\). 当 n 趋于  时, \(p^n x\) 趋于 0, 因而存在 \(n\)\(+\infty\)\(\alpha_n > 0\), 使 \(\alpha_n\) 趋于 0 且 \(\|(\phi^{-1} \circ \psi)(p^n x) - p^n x \| \leqslant \alpha_n \|p^n x\|\). 但 \(\psi(p^n x) = \psi(x)p^n\), 所以 \(\|(\phi^{-1} \circ \psi)(x) - x \| \leqslant \alpha_n \|x\|\)。)
\item
  令 U 为 A 的可逆元素的集合, 且 \(U' = 1 + m \subset U\).
\end{enumerate}

\begin{enumerate}
\def\labelenumi{(\alph{enumi})}
\item
  证明 \(\mathbf{U}' \subset \mathbf{U}_f\). (若 \(x = 1 + y\), 其中 \(y \in \mathfrak{m}\), 则由二项式定理, 当 \(x^{p^n}\)\(n\) 趋于 \(+\infty\) 时  趋于 1.)
\item
  证明 \(U_{f}\) 是 U 中像在 A/m 中为单位根的元素的集合. (使用 (a).) 从而在 K 局部紧时恢复 \(U = U_{f}\) 的事实 (此时 A/m 是有限的).
\end{enumerate}

\begin{enumerate}
\def\labelenumi{\arabic{enumi}.}
\setcounter{enumi}{3}
\tightlist
\item
  设 \(n \in \mathbf{N}^*\), p 是素数, G 是 \(p\)\(G\)\(\mathbf{GL}(n, \mathbf{Z}_p)\) 中所有元素在 p \(p\) 时模 p 同余于 1 (而在 \(p \neq 2\)\(p = 2\) 时模 4 同余于 1) 的矩阵所成的集合. 则 G 是 \(G\)\(\mathbf{GL}(n, \mathbf{Z}_p)\) 的开子群.
\end{enumerate}

\begin{enumerate}
\def\labelenumi{(\alph{enumi})}
\item
  证明 G 没有阶为 \(\neq 1\) 的有限阶元素.
\item
  证明 \(\mathbf{GL}(n, \mathbf{Z}_p)\) 的每个有限子群都同构于  时 \(\mathbf{GL}(n, \mathbf{Z} / p\mathbf{Z})\) 的一个子群 (而 \(p \neq 2\) 时同构于 \(\mathbf{GL}(n, \mathbf{Z} / 4\mathbf{Z})\) 的一个子群).\(p = 2\)
\end{enumerate}

¶ 5. (a) 设 \(\Gamma\) 是 \(G = \mathbf{GL}(n, \mathbf{Q}_p)\) 的紧子群. 证明 \(\Gamma\) 的某个共轭子群包含于 \(\mathbf{GL}(n, \mathbf{Z}_p)\). (若 T 是相对于 \(T\) 的 \(\mathbf{Q}_p^n\) 的格 (《交换代数》第 VII 章 § 4 定义 1), 证明 T 在 \(\mathbf{Z}_p\)\(T\)\(\Gamma\) 中的稳定子是 \(\Gamma\) 的开子群, 因而指数有限, 且 \(\sum_{\gamma \in I} \gamma T\) 是一个在 \(\Gamma\) 下不变的格.)

\begin{enumerate}
\def\labelenumi{(\alph{enumi})}
\setcounter{enumi}{1}
\item
  推出 \(\mathbf{G}_f\) 是 \(\mathbf{GL}(n,\mathbf{Z}_p)\) 的各共轭子群的并.
\item
  推出  的每个有限子群 \(\Phi\) 的阶都是 \(\mathbf{GL}(n, \mathbf{Q}_p)\)\(a_n(p)\) 的一个因子, 其中 \(a_n(p)\) 定义为
\end{enumerate}

\[
\begin{array}{l l} a _ {n} (p) = (p ^ {n} - 1) (p ^ {n} - p) \ldots (p ^ {n} - p ^ {n - 1}) & \text {if} p \neq 2 \\ a _ {n} (p) = 2 ^ {n ^ {2}} (2 ^ {n} - 1) (2 ^ {n} - 2) \ldots (2 ^ {n} - 2 ^ {n - 1}) & \text {if} p = 2. \end{array}
\]

(利用 (a), 可假设 \(\Phi \subset \mathbf{GL}(n, \mathbf{Z}_p)\). 然后使用练习 4.)

\begin{enumerate}
\def\labelenumi{(\alph{enumi})}
\setcounter{enumi}{3}
\tightlist
\item
  证明 \(\mathbf{SL}(n, \mathbf{Q}_p)\) 的每个有限子群的阶都是 \(s_n(p)\) 的一个因子, 其中 \(s_n(p) = \frac{a_n(p)}{p - 1}\) (当 \(p \neq 2\)) 且 \(s_n(2) = \frac{a_n(2)}{2}\).
\end{enumerate}

¶ 6. 本练习中, l 表示一个 \(\neq2\) 的素数. 若 a 是一个 \(\neq0\) 的整数, 令 \(v_{l}(a)\) 表示 a 的 l 进赋值, 即满足 \(a\equiv0\pmod{l^{e}}\) 的最大整数 e.

\begin{enumerate}
\def\labelenumi{(\alph{enumi})}
\tightlist
\item
  令 \(m\) 为一个 \(\geqslant 1\) 的整数. 记
\end{enumerate}

\[
\begin{array}{l l} \varepsilon (l, m) = 0 & \text { if } m \neq 0 (\mathrm{mod} (l - 1)) \\ \varepsilon (l, m) = v _ {l} \left(\frac {m}{l - 1}\right) + 1 & \text { if } m \equiv 0 (\mathrm{mod} (l - 1)). \end{array}
\]

证明若 x 是与 l 互素的整数, 则\(x\)\(l\)

\[
v _ {l} (x ^ {m} - 1) \geqslant \varepsilon (l, m)
\]

并且若 x 在循环群 \(x\)\((\mathbf{Z} / l^2\mathbf{Z})^*\) 中的像是该群的生成元, 则等号成立.

\begin{enumerate}
\def\labelenumi{(\alph{enumi})}
\setcounter{enumi}{1}
\tightlist
\item
  令 \(n\) 为一个 \(\geqslant 1\) 的整数. 记
\end{enumerate}

\[
r (l, n) = \sum_ {m = 1} ^ {n} \varepsilon (l, m).
\]

证明

\[
r (l, n) = \left[ \frac {n}{l - 1} \right] + \left[ \frac {n}{l (l - 1)} \right] + \left[ \frac {n}{l ^ {2} (l - 1)} \right] + \dots
\]

其中符号 \([\alpha]\) 表示实数 \(\alpha\) 的整数部分. 证明若 x 是与 l 互素的整数, 则

\[
v _ {l} ((x ^ {n} - 1) (x ^ {n - 1} - 1) \dots (x - 1)) \geqslant r (l, n)
\]

并且若 x 在 \(x\)\((\mathbf{Z} / l^2\mathbf{Z})^*\) 中的像是该群的生成元, 则等号成立.

\begin{enumerate}
\def\labelenumi{(\alph{enumi})}
\setcounter{enumi}{2}
\item
在练习 5 (c) 的记号下, 证明 \(l^{r(l,n)}\) 是当 p 为 \(l\) 的素数时整除所有 \(a_{n}(p)\) 的 l 的最大幂 (或者对所有充分大的 p, 这两种说法等价). (对 \(p\)\(\neq 1\)\(p\)\(x = p\) 应用 (b); 然后选取 p, 使它在 \(p\)\((\mathbf{Z}/l^2\mathbf{Z})^*\) 中的像是该群的生成元; 根据算术级数定理这是可能的.†)
\item
  令 \(\bar{\Gamma}\) 为 \(\mathbf{GL}(n, \mathbf{Q})\) 的有限子群, 令 \(l^e\) 为整除 \(l\)\(\Gamma\) 的阶的最大 l 幂. 证明 \(e \leqslant r(l, n)\). (利用练习 5 证明 \(l^e\) 整除所有 \(a_n(p)\), 然后应用上面的 (c).)
\item
  反之, 证明存在 \(l\) 的有限 l-子群 \(\Gamma_{i,n}\), 其阶为 \(\mathbf{GL}(n,\mathbf{Q})\)\(l^{d,l,n}\). (将其归约为 n 形如 \(n\)\(l^a (l - 1)\) 且 \(a\geqslant 0\) 的情形. 将 \(\mathbf{Q}^{\mathrm{n}}\) 分解为 \(l^a\) 个 \(\mathbf{Q}^{j - 1}\) 的副本之和, 并利用这一分解在 \(\mathbf{Q}^{\mathrm{n}}\) 上定义对称群 \(\mathrm{H}_{l,n}\)\(\mathfrak{S}_{la}\) 与群 \((\mathbf{Z} / l\mathbf{Z})^{a}\) 的半直积  的作用. 令 \(\Gamma_{l,n}\) 为 \(l\)\(\mathrm{H}_{l,n}\) 的一个 Sylow l-子群.)
\end{enumerate}

证明若 n 为偶数, 则 \(n\)\(\Gamma_{l,n}\) 包含于 \(\mathbf{Sp}(n,\mathbf{Z})\) 的某个共轭子群中.

\begin{enumerate}
\def\labelenumi{(\alph{enumi})}
\setcounter{enumi}{5}
\tightlist
\item
  *令 \(\Gamma\) 为 \(\mathbf{GL}(n, \mathbf{Q})\) 的有限子群, 且是一个 l-群. 证明 \(l\)\(\Gamma\) 共轭于 \(\Gamma_{l,n}\) 的某个子群. (首先利用适当的模 p 约化, 证明 \(p\)\(\Gamma\) 的约化是 \(\Gamma_{l,n}\) 的某个共轭子群的约化的有限子群; 然后使用 \(\Gamma\) 在 \(\mathbf{Q}^{n}\) 上的表示的特征标.) 特别地, \(\mathbf{GL}(n, \mathbf{Q})\) 中阶为 \(l^{(l,n)}\) 的每个子群都共轭于 \(\Gamma_{l,n*}\).
\end{enumerate}

¶ 7. 在本练习中, 令 \(v_{2}(a)\) 表示整数 a 的 2 进赋值.\(a\)

\begin{enumerate}
\def\labelenumi{(\alph{enumi})}
\item
  令 \(\Gamma\) 为 \(\mathbf{GL}(n,\mathbf{Q})\) 的有限子群. 证明存在一个系数为整数且正定的二次型, 在 \(\mathbf{Z}\)\(\Gamma\) 下不变. 推出 (用与练习 4 和 5 相同的论证), 对所有充分大的 p, \(p\)\(\Gamma\) 同构于域  上正交群 \(\mathbf{O}(n)\) 的一个子群 (若 \(\mathbf{F}_p\)\(\mathbf{SO}(n)\)\(\Gamma\) 包含于 \(\mathbf{SL}(n,\mathbf{Q})\), 则同构于  的一个子群).
\item
  假设 \(\Gamma\) 包含于 \(\mathbf{SL}(n, \mathbf{Q})\), 令 \(2^e\) 表示整除 \(\Gamma\) 的阶的最大 2 幂. 证明若 n 为奇数, 则 \(n\)\(2^e\) 整除所有整数
\end{enumerate}

\[
b _ {n} (p) = \left(p ^ {n - 1} - 1\right) \left(p ^ {n - 3} - 1\right) \dots \left(p ^ {2} - 1\right)
\]

(使用 (a) 和《代数》第 IX 章 § 6 的练习 13.)\(p\)

证明若 n 为偶数且 p 是充分大的素数, 则 \(n\)\(p\)\(2^e\) 整除整数\(b_n(p)\)

\[
\begin{array}{l} (p ^ {n} - 1) (p ^ {n - 2} - 1) \dots (p ^ {2} - 1) / (p ^ {n / 2} + 1) \\ (p ^ {n} - 1) (p ^ {n - 2} - 1) \dots (p ^ {2} - 1) / (p ^ {n / 2} - 1). \end{array}
\]

(与 n 为奇数时相同的方法.)\(n\)

\begin{enumerate}
\def\labelenumi{(\alph{enumi})}
\setcounter{enumi}{2}
\tightlist
\item
在 (b) 的假设下, 记
\end{enumerate}

\[
r (2, n) = n + \left[ \frac {n}{2} \right] + \left[ \frac {n}{4} \right] + \dots = n + v _ {2} (n!).
\]

令 A 为一个 \(\geqslant 3\) 的整数. 证明 \(2^{r(2,n)-1}\) 是整除所有  时的 \(b_n(p)\) 的 2 的最大幂. (与练习 6 中相同的方法†; 使用存在 \(p \geqslant A\)\(p \geqslant A\) 且 \(p \equiv 5 \pmod{8}\) 的素数这一事实.) 推出不等式 \(e \leqslant r(2,n) - 1\).

\begin{enumerate}
\def\labelenumi{(\alph{enumi})}
\setcounter{enumi}{3}
\tightlist
\item
反之, 令 \(C_{n}\) 为由置换矩阵和系数为  的对角矩阵生成的 \(\mathbf{GL}(n,\mathbf{Z})\) 的子群. \(\pm1\)\(C_{n}\) 的阶为 \(2^{n}n!\), 且 \(v_{2}(2^{n}n!) = r(2,n)\). 若 \(\Gamma_{2,n}\) 表示 \(C_{n}\) 的一个 Sylow 2-子群与 \(\mathbf{SL}(n,\mathbf{Z})\) 的交, 推出 \(\Gamma_{2,n}\) 的阶为 \(2^{n(2,n)-1}\).
\end{enumerate}

\begin{enumerate}
\def\labelenumi{\arabic{enumi}.}
\setcounter{enumi}{7}
\tightlist
\item
  令 \(n\) 为一个 \(\geqslant 1\) 的整数. 记
\end{enumerate}

\[
\mathrm{M} (n) = \prod_ {l} 1 ^ {r (1, n)},
\]

其中乘积遍历所有素数 l, 而 \(l\)\(r(l, n)\) 如练习 6 和 7 中定义.

于是 \(\mathbf{M}(1) = 2\), \(\mathbf{M}(2) = 2^3.3 = 24\), \(\mathbf{M}(3) = 2^4.3 = 48\),

\[
\mathrm{M} (4) = 2 ^ {7}. 3 ^ {2}. 5 = 5 7 6 0.
\]

由练习 6 和 7 推出, \(\mathbf{GL}(n,\mathbf{Q})\) (或 \(\mathbf{GL}(n,\mathbf{Z})\), 这两者没有区别) 的有限子群的阶的最小公倍数等于 \(\mathrm{M}(n)\). 对 \(\mathbf{SL}(n,\mathbf{Q})\) 考虑同一问题, 其中以  代替 \(\mathrm{M}(n)\).\(\frac{1}{2}\mathrm{M}(n)\)

\begin{enumerate}
\def\labelenumi{\arabic{enumi}.}
\setcounter{enumi}{8}
\item
  假设 \(\mathbf{K}\) 是局部紧的. 令 \(\mathrm{G} = \mathbf{GL}(n,\mathbf{K})\). 令 \(\mathbf{G}_1\) 为保持  的某个相对于 A 的格不变的 \(g\in \mathbf{G}\) 的集合. 令 \(\mathbf{K}^n\)\(\mathbf{G}_2\) 为在 G 中生成相对紧子群的 \(g\in \mathbf{G}\) 的集合. 令 \(\mathbf{G}_3\) 为其在 \(g\in \mathbf{G}\)\(\mathbf{K}\) 的代数闭包中的特征值绝对值为 1 的  的集合. 则 \(\mathbf{G}_1 = \mathbf{G}_2 = \mathbf{G}_3 = \mathbf{G}_f\). (使用练习 5 (a) 的论证.)
\item
  假设 K 是局部紧的. 令 G 为 K 上 n 维标准群. 令 \(\mu\) 为加法群 \(K^{n}\) 上的 Haar 测度. 证明 \(\mu|G\) 是 G 上的左、右 Haar 测度. (使用 G 是各 \(G(\alpha_{\lambda})\) 的逆极限这一事实以及《积分论》第 VII 章 §1 命题 7.)
\end{enumerate}

